\begin{textsample}{POS Dim 1 – rn – Score 31.00 – rn/rn\_em\_6.txt}  \label{ex:f1_pos_236}
2. \textbf{FUNDAMENTOS} E CONCEPÇÕES DO REFERENCIAL CURRICULAR Este Referencial Curricular apresenta -se pautado em \textbf{fundamentos} e concepções de \textbf{uma} proposta de educação democrática, integral e inclusiva para o Rio Grande do Norte, construído a partir de orientações da Base Nacional Comum Curricular - BNCC, buscando pontos de convergência possíveis e necessários.
As concepções \textbf{que} sustentam este Referencial advêm das políticas nacionais, do arcabouço epistemológico e do debate entre os membros da equipe multidisciplinar de redação curricular, entre as escolas e as instituições de ensino superior, com a interação com o Conselho Estadual de Educação.
Como orientação do projeto de educação para o Ensino Médio Potiguar, destacam -se os seguintes princípios : - Igualdade de condições para o acesso, a permanência e sucesso na escola, com inclusão, liberdade de aprender, ensinar, pesquisar e divulgar a cultura, o pensamento, a arte e o saber, atendendo à pluralidade e à diversidade cultural, resgatando e respeitando as várias manifestações de cada comunidade ; - Adoção de práticas pedagógicas \textbf{que} assegurem os direitos de aprendizagem dos estudantes, rompendo com estruturas fragmentadas do conhecimento ; - Formação integral dos \textbf{sujeitos} por meio de \textbf{um} processo de gestão da aprendizagem \textbf{ancorado} em valores éticos, políticos e estéticos ; - Valorização da experiência extraescolar e dos diferentes espaços e tempos educativos, abrangendo espaços sociais na escola e fora dela ; - Vinculação entre a educação escolar, as práticas sociais e o mundo do trabalho ; - Gestão democrática do ensino público, na forma da legislação e das normas dos respectivos sistemas de ensino.
Considera -se importante mencionar os avanços trazidos pela Emenda Constitucional nº 59/2009 e pela Lei de Diretrizes e Bases da Educação Nacional nº 9.394/1996, \textbf{que} consolidaram a obrigatoriedade da educação básica gratuita dos 4 ( quatro ) aos 17 ( dezessete ) anos de idade, igualmente, assegurando a oferta gratuita e acesso para aqueles \textbf{que} não tiveram na idade própria.
Assim, o Referencial Curricular do Rio Grande do Norte busca contribuir com subsídios \textbf{inerentes} à transformação premente e necessária da práxis da educação básica no Estado, ofertando educação igualitária e equitativa, de forma a promover mudanças na sociedade potiguar \textbf{que} contribuam para a superação da desigualdade social marcante na sociedade brasileira.
Elegemos princípios e práticas para nortearem o trabalho do conjunto das unidades escolares, incluindo sugestões didáticas \textbf{que} orientam e dinamizam o planejamento e a realização do trabalho pedagógico com os estudantes, pensados nas áreas e componentes curriculares do Ensino Médio \textbf{que} atendam suas modalidades e ofertas específicas.
O panorama \textbf{atual} do ensino médio \textbf{exigiu} \textbf{que} as unidades escolares se reinventassem para conceber as diversas mudanças na maneira de ensinar e de aprender, em especial com o advento das tecnologias digitais intensificado pelo período da pandemia do coronavírus ( COVID-19 ).
Nesse período, \textbf{que} coincidiu com a elaboração deste documento, compreende -se \textbf{que} não há como ter \textbf{uma} educação \textbf{justa}, democrática e solidária se não houver a garantia de igualdade e \textbf{equidade} no acesso ao ensino-aprendizagem para todos. \textbf{Logo}, precisamos orientar a ação \textbf{didático-pedagógica} por meio do planejamento integrado, de modo a contemplar aos estudantes em sua inteireza, visando mitigar os efeitos das desigualdades estabelecidas nas escolas.
O Sistema de Ensino Público materializa o \textbf{compromisso} de inclusão, promovendo apoio às escolas e buscando condições viáveis à efetivação da aprendizagem de todos os seus estudantes, numa perspectiva de garantir as mesmas oportunidades de aprender, embora utilizando \textbf{uma} diversidade de estratégias para esse ponto de chegada.
Toda essa diversidade \textbf{exige} práticas pedagógicas \textbf{que} dialoguem com a inclusão e \textbf{um} currículo \textbf{que} traga a voz dos silenciados e excluídos, escrevendo \textbf{uma} nova história sem ferir o direito à aprendizagem, sob os princípios da igualdade e da \textbf{equidade}.
Para tanto, recorremos à Lei estadual 10.659/19, denominada de “ Lei da Escola Democrática ”, \textbf{que} no seu Art. 1º institui alguns princípios \textbf{inerentes} à educação democrática do Estado do Rio Grande do Norte : > Art. 1º Todos os professores, estudantes e funcionários são livres para expressar seus pensamentos e suas opiniões no ambiente escolar das redes pública e privada de ensino do Rio Grande do Norte, em consonância com os seguintes princípios : > I - liberdade de aprender, ensinar, pesquisar e divulgar o pensamento, a arte e o saber ; > II - pluralismo de ideias e de concepções pedagógicas ; > III - respeito à liberdade e apreço à tolerância ; > IV - \textbf{ideais} de solidariedade humana para o pleno desenvolvimento do \textbf{educando} ; > V - preparo para o exercício da cidadania e sua qualificação para o trabalho.
O texto da referida lei descreve o espaço escolar como \textbf{um} ambiente democrático, aberto a diferentes estratégias de diálogos e \textbf{exposições} de visões críticas, aspectos esses intrínsecos à construção de \textbf{uma} sociedade \textbf{justa}, garantindo aos profissionais da educação a liberdade de \textbf{desempenharem} o magistério longe de qualquer indício de ações repressivas e/ou censura.
De acordo com Silva ( 1999 ), no currículo forja -se a identidade, o \textbf{que} \textbf{implica} a ênfase de questões \textbf{identitárias} como ponto de atenção para a execução do trabalho com as juventudes na condição de \textbf{sujeitos} do ensino médio.
Este Referencial Curricular não tem a pretensão de ser \textbf{um} guia ou \textbf{uma} ferramenta de modelagem do trabalho pedagógico, mas assumir o papel de \textbf{norteador}, condutor e orientador da instituição de ensino em sua organização e funcionamento, \textbf{baseando} -se nas potencialidades, aspirações e necessidades da comunidade em questão.
Segundo Silva e Carvalho ( 2019 ), a ideia de educação integral está relacionada à formação de indivíduos capazes de compreender e modificar o meio em \textbf{que} \textbf{vivem}, buscando sempre o bem comum e a convivência solidária com os seus pares.
Neste Referencial, a educação integral visa à formação e ao desenvolvimento humano em todos os seus aspectos, o \textbf{que} requer o entendimento da complexidade e a não linearidade desse desenvolvimento, \textbf{revelando} \textbf{uma} visão plural, singular e integral do adolescente, do jovem e do adulto, considerando -os como \textbf{sujeitos} de aprendizagem respeitados em suas \textbf{singularidades} e diversidades.
A formação e o desenvolvimento \textbf{humanos} são \textbf{imprescindíveis} sob \textbf{uma} visão global e, ao mesmo tempo, singular, quanto aos aspectos biopsicossociais e afetivos dos \textbf{sujeitos}.
Assim, a escola deve oportunizar aos estudantes \textbf{um} ambiente propício ao atendimento de sua formação, de modo \textbf{que} promova o protagonismo dos jovens, \textbf{enquanto} sujeitos críticos, emancipados e cooperativos.
Nesse sentido, serão capazes de atuar na sociedade contribuindo para a construção e assunção de \textbf{um} discurso de alteridade, bem como de \textbf{uma} práxis \textbf{que} considera a diversidade, com disponibilidade para resolução de problemas e consequente tomada de decisões.
Quando nos referimos à educação inclusiva, pensamos na garantia do acesso à educação de qualidade a todos, respeitando as diferenças e as diversidades.
Para tanto, é essencial exercitar o princípio da \textbf{equidade} por meio de estratégias \textbf{que} evidenciem as particularidades dos \textbf{sujeitos}, considerando suas necessidades e características para \textbf{que} todos atinjam o mesmo ponto de chegada – a aprendizagem.
Portanto, não podemos deixar de \textbf{ver} a importância político-social do currículo no estabelecimento de \textbf{uma} prática emancipatória e inclusiva na escola, em detrimento da \textbf{exclusão}, a fim de \textbf{que} os processos de constituição da cultura, seus conflitos e contradições possam resgatar vozes silenciadas no território escolar ( Bezerra e Ribeiro 2010 ).
Assim, confirmamos a assunção da visão inclusiva do Referencial Curricular do Ensino Médio Potiguar, quando \textbf{preconiza} o respeito à individualidade, o direito à acessibilidade da pessoa com deficiência, com transtornos globais do desenvolvimento e daqueles com altas habilidades ou superdotação.
Igualmente, propomos o acolhimento e a inclusão pelas escolas potiguares de estudantes de diversas etnias, culturas, gêneros, orientação sexual, territórios, diferentes expressões, entre outros.

% matched lemmas: ancorar, atual, basear, compromisso, desempenhar, didático-pedagógico, educar, enquanto, equidade, exclusão, exigir, exposição, fundamento, humanos, ideal, identitária, implicar, imprescindível, inerente, justo, logo, norteador, preconizar, que, revelar, singularidade, sujeito, um, ver, viver
\end{textsample}
